\documentclass[11pt,letterpaper]{article}
\usepackage[T1]{fontenc}
\usepackage{lmodern,amsmath,amssymb,amsthm,microtype}
\usepackage[margin=1in]{geometry}
\usepackage[hidelinks]{hyperref}
\providecommand{\tightlist}{\setlength{\itemsep}{0pt}\setlength{\parskip}{0pt}}
\DeclareUnicodeCharacter{2264}{\ensuremath{\le}}
\DeclareUnicodeCharacter{2265}{\ensuremath{\ge}}
\DeclareUnicodeCharacter{2260}{\ensuremath{\ne}}
\DeclareUnicodeCharacter{2212}{\ensuremath{-}}

\DeclareUnicodeCharacter{2211}{\ensuremath{\sum}}

\usepackage{longtable,booktabs,array}
\title{An anchored-tail exclusion for nine separated subset sums below 150}
\author{Aafi Khan Malek\\ \small Independent researcher, Ahmedabad, India\\
\small\href{mailto:aafimalek2032@gmail.com}{aafimalek2032@gmail.com}}
\date{28 November 2026}
\begin{document}
\maketitle
\textbf{Theorem.} Let nine sorted positive real weights have pairwise
unit-separated subset sums. If \(149\le x_9<150\), then \(x_8\ge133\).
We use the previously proved
\href{https://github.com/Aafimalek/novel-discoveries/tree/1fdefaee1c13f9315b29c352b7bba3231b48585d/nine-separated-subset-sums/verification/rank-cuts-below-150/penultimate-cap-131}{cap-131 exclusion}
\(x_8\ge132\), which bounds the last ordered deficit by one in the
next cap. The remaining finite exclusions use exact integer arithmetic
and are not Lean-verified. The largest-weight interval is not closed.

\section{Cap-specific inequalities}\label{cap-specific-inequalities}

For cap \(t=132\) the forced coordinate-9 starts are 143 through 172, so
\(q=35,\dots,64\). The common counts are \(64\le a\le74\) and
\(b=148-a\), and \(\alpha_q=u_{a-q+1}\in[149-2q,151-2q)\) are
consecutive base unsigned full sign magnitudes. With the earlier
plus-two endpoints, the capacity ranges for all 128 complementary signs
are:

\begin{longtable}[]{@{}lr@{}}
\toprule\noalign{}
Half-open magnitude range & Capacity \\
\midrule\noalign{}
\endhead
\bottomrule\noalign{}
\endlastfoot
{[}1,21) & \(a-64\) \\
each {[}149−2q,151−2q), q from 64 down to 35 & 1 \\
{[}81,219) & 34 \\
{[}278,321) & \(148-2a\) \\
{[}379,10000) & \(a-20\) \\
\end{longtable}

The two holes have width at least \(2(64-35)+4=62\). A flip at weight
strictly below 31 or a swap at adjacent weight gap strictly below 31
preserves each broad group with rank indices in {[}0,a), {[}a,b), or
{[}b,128); these are zero-based position sets after expanding each bin
by its capacity. Empty bins, including the middle broad group if a=b,
add no positions. Small flip orbits give \(x_3\ge31\) since
\(2^s\mid a,b\) would imply \(2^s\mid148\) and \(s\le2\).

For \(y=(126,\dots,133)\) write \(r_i=y_i-x_i\) and \(R=\sum r_i\).
Unbalanced signed magnitudes at y are at least 244 and balanced ones at
most 16. If any prescribed sign is unbalanced, then \(R>244-81=163\).
Otherwise the balanced \(\alpha_{35}\ge79\) gives \(R\ge63\). Hence in
all cases \(R\ge63\).

\section{The anchored tail and finite worker
proof}\label{the-anchored-tail-and-finite-worker-proof}

The previous cap proof gives \(x_8\ge132\), so the last deficit is
\(z=r_8\le1\). For any ordered tail of N deficits with total T, subtract
z from all entries, making the last entry zero. By prefix differences,
the signed contribution of the first N−1 residual entries lies between
their extreme prefix-sign averages \(\mu\) times \(T-Nz\). After adding
back the common contribution, the bounding expression is
\(\mu T+(\sum\sigma-N\mu)z\). Use the tail interval for T and
\(0\le z\le\min(1,T_{\rm upper}/N)\) to bound the terms independently,
rounding outward. This is a safe relaxation even if the two extrema
cannot occur together.

The two-deficit stage starts with 7,875 unit cells; 6,966 reach the
count splits. Across 76,626 count profiles, there are 38,640 pure
orbit-size rejections, 2,610 individual interval rejections, and 26,950
interval-conditioned orbit rejections, leaving \textbf{8,426}. The
ordinary interval and sign-orbit packing arguments are the same as in
cap 131. Refine each of the 8,426 cases by the third deficit (and, only
when needed, the fourth), retaining the anchored signed bounds.

A deterministic phase-2 replay generated the exact residual profile
payload, which is split into \textbf{70 disjoint exhaustive shards} of
at most 128 inputs each. All 70 complete with zero residual: 67 at depth
3 and three at depth 4. In aggregate:
\begin{itemize}
\item Depth 3 covers 387,410 cells:
123,021 orbit-size rejections, 62,657 ordinary interval rejections,
201,702 interval-conditioned orbit rejections, and 30 residual cells.
\item Depth 4 covers 952 refinements of those 30 cells: 818, 52, and 82 by the
respective three tests, with none remaining.

\end{itemize}

The earlier monolithic run timed out and is not used for third-stage
coverage. Each accepted shard has a separately retained log and an
input/output SHA chain from the exact phase-2 payload. One worker ran at
a time, with 35-second and 143,000-KiB virtual-memory guards per
invocation.

\section{Independent audit and
replay}\label{independent-audit-and-replay}

\begin{verbatim}
sha256sum -c SHA256SUMS
python3 coverage_audit.py
./replay-all.sh
\end{verbatim}

The first two commands independently check source/payload/manifest
integrity, disjoint full shard coverage, and that all stored chains end
empty. For a fresh exact computation, the last command first reruns
phase 2 and compares its canonical byte payload; then it serially
recomputes all 70 shards to a fresh results directory. An individual
shard can instead be run with a separate results directory using the
in-packet run-shard.sh wrapper. Nothing in the verifier depends on
third-party solver output or floating arithmetic. The mathematical rank
and endpoint implications are handwritten, not a Lean theorem.


\end{document}
